\chapter{Literature and Background}

\section{Evolutionary design and variable-length genotypes}

Evolutionary design is the use of computational search to automatically generate or improve artefacts whose quality can be assessed by an objective function. In contrast to conventional optimization problems with a small number of predefined numerical parameters, evolutionary design often concerns structured objects: shapes, mechanisms, controllers, artificial organisms, or other constructions whose form is not known in advance. The search process therefore has to explore both parameter values and structural alternatives. Candidate solutions are usually represented as genotypes, transformed into phenotypes, and evaluated according to a task-dependent fitness criterion. Selection and variation operators then bias the population toward solutions with better observed performance. % TODO: add citation for evolutionary design and evolutionary computation.

A distinctive feature of many evolutionary design problems is that the genotype does not have a fixed and directly positional interpretation. In a fixed-length vector representation, the meaning of a coordinate is typically stable across all individuals: the same index denotes the same parameter. By contrast, a variable-length genotype can contain fragments whose role depends on their position, nesting, and surrounding symbols. Such genotypes may encode branching morphology, repeated substructures, hierarchical components, or program-like construction rules. Their length may change during mutation and recombination, and a local edit can alter not only one parameter but also the interpretation of a larger downstream structure.

This property makes evolutionary design more expressive, but it also makes optimization more difficult. The search space becomes discrete, irregular, and strongly constrained by syntactic and semantic validity. Some candidate strings may not describe any admissible phenotype, while others may be valid but differ greatly in behaviour after only a small textual change. Furthermore, the contribution of a genotype fragment is often context-dependent: a fragment that is useful in one structural location can be neutral, invalid, or harmful in another. This context dependence is closely related to epistasis, understood broadly as interaction between parts of the representation, where the effect of modifying one part depends on other parts of the genotype. % TODO: add citation for epistasis and structured evolutionary representations.

Simulation-based optimization adds another layer of complexity. In this setting, fitness is not obtained from a closed-form mathematical expression, but from a simulation that constructs and evaluates the phenotype. This is typical for artificial-life and evolutionary robotics environments, where the performance of a morphology may depend on physical interaction, dynamics, control, or environmental conditions. Simulation enables the assessment of complex behaviours, but it also makes each evaluation relatively expensive and sometimes noisy. As a result, variation operators should ideally produce candidates that are not only diverse, but also valid and promising enough to justify evaluation.

The representation used in this thesis belongs to this class of problems. Framsticks F1 genotypes are variable-length symbolic descriptions of virtual organisms, and their fitness is obtained after phenotype construction and simulation. The \texttt{vertpos} criterion considered later in the thesis is therefore an evaluated property of the resulting organism, not a simple property of the raw character string. This distinction is important because a method that performs well at modifying genotype text does not necessarily perform well at improving simulated fitness. Conversely, an optimization method may be valuable if it discovers modifications that preserve validity and produce useful phenotypic effects, even when the mapping from genotype to fitness is highly indirect.

The background developed in this section motivates the two representation-aware directions studied in the thesis. Compatible Substitutions Optimization (CoSO) addresses the difficulty directly in the discrete genotype representation by considering compatibility between a substitution and its context. Latent Space Optimization (LSO) approaches the same difficulty by learning a continuous representation in which variable-length genotypes are mapped to fixed-dimensional vectors before search is performed. Both directions can be understood as attempts to define more useful neighbourhoods for structured genotypes than those obtained from naive position-independent edits. The following sections introduce the concrete Framsticks F1 representation and then discuss these two approaches in greater detail.

\section{Framsticks and the F1 genetic representation}

Framsticks is an artificial-life and evolutionary-design environment in which virtual organisms are described by genotypes, developed into phenotypes, simulated, and evaluated according to selected criteria. % TODO: add citation for Framsticks.
For the purposes of this thesis, Framsticks provides both the design domain and the final evaluator of candidate solutions. The methods compared later operate on, or generate, symbolic genotype descriptions, but their usefulness is ultimately determined by whether these descriptions can be interpreted as valid organisms and assessed in simulation. The environment therefore connects representation-level search with phenotype-level evaluation.

The representation considered in this thesis is Framsticks F1. At a conceptual level, an F1 genotype is a textual description of an organism rather than a fixed list of numerical parameters. Its symbols specify structural components and modifiers, and the complete string is interpreted through a genotype-to-phenotype mapping. This mapping is indirect: the fitness of a candidate is not defined by the string alone, but by the organism produced from it and by the behaviour of this organism during simulation. The later use of the \texttt{vertpos} criterion should therefore be understood as evaluation of a constructed phenotype, not as a direct score assigned to individual characters.

F1 genotypes are variable-length expressions. This property is important because evolutionary variation can change not only parameter values, but also the amount and organization of encoded structure. Longer or shorter genotypes may describe organisms with different numbers of components, different branching patterns, or different modifier sequences. Consequently, the search space is not a simple Euclidean space with a fixed dimensionality. It is a discrete structured space in which insertion, deletion, and substitution can change the interpretation of neighbouring fragments and, in some cases, larger parts of the organism.

The F1 representation also contains branching and nested descriptions. In simplified grammatical terms, genotype strings may be understood as sequences of elements, where elements can contain a core component, optional modifiers, and branch expressions enclosed in parentheses. Branches make the representation expressive, because they allow the genotype to describe articulated or tree-like morphology. At the same time, they impose syntactic constraints: parentheses, separators, branch contents, and modifier positions must be arranged in forms that can be parsed and developed into a phenotype. A candidate that violates these constraints may be useless for simulation-based optimization because it cannot be evaluated in the intended way.

Syntactic validity is therefore a central issue for both approaches studied in the thesis. In a direct genotype-level method, a local edit must preserve enough structure for the resulting expression to remain acceptable. In an autoencoder-based method, a decoder must produce strings that are not merely character sequences, but valid F1 genotypes. Validity is not equivalent to good performance, but invalidity prevents meaningful phenotype evaluation. This distinction motivates the later separation between reconstruction or validity analysis and actual optimization performance.

Another important property of F1 is context sensitivity. The effect of a symbol or fragment depends on where it appears in the genotype and on the structural role of the surrounding expression. A modifier before a core component, a symbol inside a branch, and a separator between branches are not interchangeable even if they are all local textual tokens. Therefore, edits that ignore context may easily produce invalid genotypes or changes whose phenotypic consequences are difficult to control. This observation prepares the ground for context-aware genotype modification and for Compatible Substitutions Optimization, where the compatibility of a substitution with its surrounding expression is treated as an essential concern.

The same representational properties also motivate Latent Space Optimization. Since F1 genotypes are variable-length, branching, and syntactically constrained, mapping them to fixed-dimensional continuous vectors is non-trivial. A useful latent representation should preserve enough information about the genotype structure to allow decoding back into valid candidates, while also arranging genotypes in a way that supports search. The background role of the F1 representation is thus twofold: it defines the concrete optimization domain of the thesis, and it explains why both explicit context-aware edits and learned continuous representations are relevant alternatives to naive character-level modification.

\section{Context-aware genotype modification}

% TODO: write this section.

\section{Compatible Substitutions Optimization}

% TODO: write this section.

\section{Autoencoders and variational autoencoders}

Autoencoders are representation-learning models trained to reproduce their inputs after passing them through an intermediate code. The model is usually divided into an encoder and a decoder. The encoder maps an input object to a latent representation, while the decoder attempts to reconstruct the original object from this representation. If the latent representation is constrained to be lower-dimensional, regularized, or otherwise information-limited, the model is encouraged to capture regularities that are useful for reconstruction rather than simply copying the input \cite{Rumelhart1986,HintonSalakhutdinov2006}.

This encoder--decoder structure is particularly relevant for the study of structured genotypes. A genotype such as a Framsticks F1 string is variable-length and symbolic, whereas many optimization methods operate most naturally on fixed-size numerical vectors. An autoencoder can therefore be used as a bridge between these two forms. The encoder converts a genotype into a vector in the latent space, and the decoder converts a latent vector back into a candidate genotype. In this thesis, this idea is important not because reconstruction is an end in itself, but because reconstruction defines whether the learned representation can be used as a search space for genotype generation and modification.

Variational autoencoders extend this idea by making the latent representation probabilistic. Instead of encoding an input as a single unconstrained vector only, a variational autoencoder represents the input by parameters of a latent distribution, typically mean and variance parameters, and samples a latent vector from this distribution during training. The training objective combines reconstruction accuracy with a regularization term that encourages the latent distribution to remain close to a simple prior, often a standard normal distribution. This regularization is intended to make the latent space more continuous and more suitable for sampling, interpolation, and optimization than an arbitrary collection of isolated codes \cite{KingmaWelling2014,RezendeMohamedWierstra2014}.

For continuous data, reconstruction can often be measured by numerical distance. For discrete structured objects, however, reconstruction quality has additional dimensions. A decoded sequence may differ from the input by only a small number of symbols and still be syntactically invalid. Conversely, a reconstruction that is textually different from the input may still describe a valid and potentially useful organism. Therefore, when autoencoders are applied to F1 genotypes, the central concerns are not limited to token-level accuracy. The decoded output must also respect the constraints of the representation so that it can be interpreted as a valid Framsticks genotype and evaluated with a criterion such as \texttt{vertpos}.

This validity issue explains why grammar- and structure-aware decoders are relevant. A purely character-level decoder treats generation mainly as a sequence prediction task and must learn syntactic constraints indirectly from examples. This may be sufficient for simple patterns, but it provides no explicit guarantee that parentheses, branches, modifiers, and core elements will appear in admissible combinations. By contrast, grammar-aware representations can decode sequences of production rules rather than raw characters, and masked variants can restrict choices to rules that are valid for the currently expanded nonterminal \cite{KusnerPaigeHernandezLobato2017}. Tree-oriented variants can additionally reflect the fact that F1 genotypes contain branching structure rather than only a flat string. Such mechanisms do not by themselves guarantee high fitness, but they can reduce the burden placed on the model by embedding part of the known syntax into the decoder.

In the context of this thesis, autoencoders and variational autoencoders therefore provide the background for the Latent Space Optimization part of the comparison. They define a learned continuous representation of F1 genotypes that is different from the explicit discrete representation used by context-aware substitutions and CoSO. The usefulness of this representation depends on several separate properties: whether inputs can be encoded and reconstructed, whether decoded samples remain valid genotypes, whether nearby latent points correspond to related genotype structures, and whether movement in the latent space is correlated with changes in simulation-based fitness. These properties are connected, but none of them alone is sufficient to establish that an autoencoder is useful for optimization.

\section{Latent Space Optimization}

Latent Space Optimization (LSO) is the use of a learned latent representation as the domain in which search is performed. Instead of applying variation directly to the original object, an object is encoded into a continuous vector, optimization modifies or samples vectors in this latent space, and each candidate vector is decoded back into the original representation before evaluation. Related work has used learned continuous representations for molecular design, program synthesis, and evolutionary search in learned representations, illustrating the general idea that an autoencoder or another generative model can define a search space for structured objects \cite{GomezBombarelli2018,KaszubaKomosinskiMensfelt2021,LiskowskiKrawiecTokluSwan2020,BentleyLimGaierTran2022}. In the setting considered in this thesis, the original objects are Framsticks F1 genotypes, and the final objective remains the true simulation-based \texttt{vertpos} fitness. The latent vector is therefore an intermediate representation for proposing candidates, not a replacement for evaluation in Framsticks.

The motivation for LSO is that a learned continuous representation may provide a more compact or smoother search domain than the raw genotype space. F1 genotypes are variable-length symbolic expressions with branching syntax and context-dependent effects. Direct edits in such a space can easily produce invalid candidates or large unpredictable phenotypic changes. If an autoencoder has learned useful regularities of the genotype distribution, nearby latent vectors may decode to related genotypes, and continuous perturbations may correspond to moderate structural changes. This can make it possible to apply optimization techniques that are difficult to define directly on raw F1 strings.

The general LSO loop consists of four conceptual steps. First, one or more starting genotypes are encoded into latent vectors. Second, an optimizer proposes new latent vectors by mutation, recombination, distribution update, or another continuous-space operation. Third, each proposed vector is decoded into an F1 genotype. Fourth, the decoded genotype is checked for validity and evaluated in Framsticks, where its \texttt{vertpos} value provides feedback to the optimizer. This loop preserves a common final evaluation criterion with direct evolutionary search: regardless of how a candidate is generated, it is useful only if it decodes into a genotype that can be evaluated and achieves good simulated performance.

Several optimizer families can be used within this framework. An evolutionary search in latent space keeps a population of latent vectors, selects better candidates according to decoded fitness, and creates new vectors through latent mutation, interpolation-like crossover, or related variation operators. The Cross-Entropy Method (CEM) instead maintains a sampling distribution over the latent space, evaluates a population of samples, and updates the distribution toward the best-performing elite samples \cite{Rubinstein1999}. Covariance Matrix Adaptation Evolution Strategy (CMA-ES) is another population-based continuous optimizer that adapts both the mean and covariance of its search distribution \cite{Hansen2016}. These methods differ in how they use fitness feedback, but in LSO they share the same requirement: every promising latent candidate must ultimately be decoded into a valid genotype and evaluated by the original objective.

LSO may help with structured genotypes when the latent representation captures properties that are relevant to both syntax and fitness. A compact latent space can reduce the apparent dimensionality of the search, and a regularized decoder can bias generated candidates toward patterns observed in valid genotypes. This can be advantageous when direct random edits are too disruptive. Moreover, continuous search operators can explore directions, mixtures, and local neighbourhoods in a way that has no direct analogue in symbolic genotype strings.

However, these advantages are conditional rather than guaranteed. If the autoencoder reconstructs poorly, many optimized latent vectors may decode into invalid or low-quality genotypes. If the latent space is syntactically valid but not fitness-aware, the optimizer may efficiently explore regions that reproduce common genotype patterns without improving \texttt{vertpos}. If the decoder maps large regions of latent space to similar genotypes, search may stagnate because many candidate vectors produce duplicates or near-duplicates. Conversely, if small latent changes decode into unrelated structures, the expected smoothness of the latent space is lost. For this reason, reconstruction quality, validity, diversity, and optimization performance must be analysed separately.

\section{Summary}

% TODO: write this section.
